\documentclass[11pt,a4paper]{article}

\usepackage[margin=1in]{geometry}
\usepackage{hyperref}
\usepackage{enumitem}
\usepackage{booktabs}
\usepackage{tabularx}
\usepackage{xcolor}
\usepackage{mdframed}
\usepackage{longtable}
\usepackage{array}

\hypersetup{colorlinks=true, linkcolor=blue!70!black, urlcolor=blue!70!black}

\title{Indoor Navigation Stack Breakdown\\[0.3em]
\large Exact module responsibilities, interfaces, and work allocation}
\author{Project Working Draft}
\date{\today}

\begin{document}

\maketitle
\tableofcontents
\newpage

\section{Why This Document Exists}

This document is the \textbf{operational version} of the indoor strategy.
It answers:
\begin{itemize}[nosep]
  \item what each stack is supposed to do,
  \item what goes into each stack,
  \item what comes out of each stack,
  \item what technical methods we plan to use first,
  \item who should own which stack,
  \item and how the stacks connect into one system.
\end{itemize}

\begin{mdframed}
\textbf{System goal:} indoor image-goal navigation on EarthRover using a
topological graph for global navigation, MBRA for local image-goal control, and
a safety layer for collision prevention.
\end{mdframed}


\section{System in One Pass}

\subsection*{End-to-end runtime loop}

\begin{enumerate}[nosep]
  \item read the latest front camera frame from the EarthRover SDK,
  \item update the recent image buffer,
  \item localize the current view in the indoor topological graph,
  \item map the mission checkpoint image to a target graph node,
  \item plan a graph path from current node to target node,
  \item choose the next nearby node on that path as the current subgoal image,
  \item run MBRA on recent context plus the subgoal image,
  \item run safety checks on the proposed control,
  \item send the safe control command to the robot,
  \item repeat until checkpoint image verification says the target is reached.
\end{enumerate}

\subsection*{Main stacks}

\begin{enumerate}[nosep]
  \item SDK / robot interface
  \item Perception
  \item Map / graph
  \item Planning
  \item Control
  \item Safety
  \item Integration / orchestration
  \item Testing / evaluation
\end{enumerate}


\section{Ownership Model}

If people need to be assigned immediately, use the following work split.
Replace the owner placeholders with actual names.

\begin{tabularx}{\textwidth}{p{0.18\textwidth} p{0.22\textwidth} X}
\toprule
\textbf{Stack} & \textbf{Owner} & \textbf{Primary responsibility} \\
\midrule
SDK / Interface & Owner A & camera ingestion, control send, runtime I/O reliability \\
Perception & Owner B & visual localization, checkpoint verification, pedestrian detection \\
Map / Graph & Owner C & graph data model, graph builder, node/edge storage \\
Planning & Owner D & graph search, subgoal selection, progress logic \\
Control & Owner E & MBRA runtime wrapper and local control loop \\
Safety & Owner F & depth / detector safety overrides and stop logic \\
Integration & Owner G & connect all modules into one process or service \\
Testing & Owner H & route tests, metrics, failure analysis, logging \\
\bottomrule
\end{tabularx}


\section{Stack 0: SDK / Robot Interface}

\subsection{Purpose}

Provide reliable robot I/O for the rest of the system.

\subsection{What this stack receives}

\begin{itemize}[nosep]
  \item browser-fed EarthRover camera stream,
  \item EarthRover telemetry,
  \item command requests from the navigation system.
\end{itemize}

\subsection{What this stack outputs}

\begin{itemize}[nosep]
  \item latest front camera frame,
  \item timestamps and any usable telemetry,
  \item confirmed command send path to the robot.
\end{itemize}

\subsection{What we are going to do}

\begin{enumerate}[nosep]
  \item use the existing EarthRover SDK server,
  \item use the existing Python interface to fetch frames and send commands,
  \item add any missing timestamping or logging needed for indoor experiments,
  \item standardize a clean API for the rest of the indoor pipeline.
\end{enumerate}

\subsection{Implementation details}

\textbf{Primary functions needed:}
\begin{itemize}[nosep]
  \item \texttt{get\_camera\_frame()}
  \item \texttt{send\_control(linear, angular)}
  \item optional frame timestamp retrieval
  \item optional telemetry logging
\end{itemize}

\subsection{Definition of done}

\begin{itemize}[nosep]
  \item frame retrieval is stable,
  \item command sending is stable,
  \item all indoor modules can consume the same camera and control API.
\end{itemize}


\section{Stack 1: Perception}

\subsection{Purpose}

Understand what the robot is currently seeing.

\subsection{Submodules}

\begin{enumerate}[nosep]
  \item \textbf{Context buffer}: keep the last $N$ frames for MBRA.
  \item \textbf{Visual localization}: match the current view to graph nodes.
  \item \textbf{Checkpoint verification}: determine whether a goal image is reached.
  \item \textbf{Pedestrian detection}: detect people for safety support.
\end{enumerate}

\subsection{Inputs}

\begin{itemize}[nosep]
  \item latest front camera frame from SDK,
  \item graph node images from the map stack,
  \item target checkpoint image from mission definition.
\end{itemize}

\subsection{Outputs}

\begin{itemize}[nosep]
  \item current node estimate,
  \item top-$k$ candidate matching nodes,
  \item localization confidence,
  \item checkpoint reached / not reached,
  \item pedestrian present / not present in danger zone.
\end{itemize}

\subsection{What we are going to do first}

\begin{enumerate}[nosep]
  \item implement a clean frame buffer,
  \item choose a first visual descriptor baseline for node retrieval,
  \item implement a first checkpoint verification baseline,
  \item run a pretrained person detector for corridor safety.
\end{enumerate}

\subsection{Preferred first methods}

\textbf{Visual localization:}
\begin{itemize}[nosep]
  \item start with global descriptor retrieval against graph nodes,
  \item add temporal smoothing,
  \item add local geometric verification for strong matches.
\end{itemize}

\textbf{Checkpoint verification:}
\begin{itemize}[nosep]
  \item stage 1: descriptor similarity threshold,
  \item stage 2: feature matching + geometric verification,
  \item require multiple consecutive confirmations if needed.
\end{itemize}

\textbf{Pedestrian detection:}
\begin{itemize}[nosep]
  \item start with pretrained YOLO \texttt{person} detection,
  \item define a forward danger zone in image coordinates,
  \item pass danger flags to the safety stack.
\end{itemize}

\subsection{Definition of done}

\begin{itemize}[nosep]
  \item current node can be estimated from live camera frames,
  \item checkpoint verification can confirm a reached goal image,
  \item pedestrian detector works well enough to support indoor safety.
\end{itemize}


\section{Stack 2: Map / Graph}

\subsection{Purpose}

Represent the indoor environment as a place graph instead of a GPS map.

\subsection{Inputs}

\begin{itemize}[nosep]
  \item teleoperated indoor camera recordings,
  \item route ordering,
  \item manually or semi-automatically selected keyframes.
\end{itemize}

\subsection{Outputs}

\begin{itemize}[nosep]
  \item graph node set,
  \item graph edge set,
  \item node images and metadata,
  \item target checkpoint-to-node mapping.
\end{itemize}

\subsection{What we are going to do}

\begin{enumerate}[nosep]
  \item teleoperate the robot through indoor routes,
  \item save reference images,
  \item choose representative keyframes as nodes,
  \item define traversable connections as edges,
  \item store the graph in a format the planner can query directly.
\end{enumerate}

\subsection{Minimum graph schema}

Each node should contain:
\begin{itemize}[nosep]
  \item node id,
  \item image path,
  \item route id,
  \item ordering index,
  \item optional semantic label,
  \item optional descriptor embedding.
\end{itemize}

Each edge should contain:
\begin{itemize}[nosep]
  \item source node id,
  \item destination node id,
  \item traversal cost,
  \item optional direction label.
\end{itemize}

\subsection{Definition of done}

\begin{itemize}[nosep]
  \item one indoor route can be represented as a queryable graph,
  \item the planner can load it and search it,
  \item the perception stack can localize against it.
\end{itemize}


\section{Stack 3: Planning}

\subsection{Purpose}

Decide what place the robot should go to next.

\subsection{Inputs}

\begin{itemize}[nosep]
  \item current node estimate from perception,
  \item target node from mission checkpoint mapping,
  \item graph from map stack.
\end{itemize}

\subsection{Outputs}

\begin{itemize}[nosep]
  \item graph path from current node to target node,
  \item next subgoal node,
  \item subgoal image for MBRA,
  \item path progress state.
\end{itemize}

\subsection{What we are going to do}

\begin{enumerate}[nosep]
  \item compute shortest path in the graph,
  \item pick the next nearby node on that path,
  \item expose the next node image as the current local goal,
  \item update path state as localization changes,
  \item add simple recovery if localization confidence drops.
\end{enumerate}

\subsection{Important design rule}

\begin{mdframed}
\textbf{MBRA should receive a nearby subgoal image, not the final faraway goal
image for the whole mission.}
\end{mdframed}

\subsection{Definition of done}

\begin{itemize}[nosep]
  \item planner can produce a valid path,
  \item planner can switch subgoals correctly,
  \item control stack can always request the current subgoal image.
\end{itemize}


\section{Stack 4: Control}

\subsection{Purpose}

Move the robot from the current view toward the next nearby subgoal image.

\subsection{Inputs}

\begin{itemize}[nosep]
  \item recent frame buffer from perception,
  \item current subgoal image from planner,
  \item optional delay estimate,
  \item recent command history,
  \item robot-size inference constant if required by MBRA.
\end{itemize}

\subsection{Outputs}

\begin{itemize}[nosep]
  \item proposed linear command,
  \item proposed angular command,
  \item controller diagnostics.
\end{itemize}

\subsection{What we are going to do}

\begin{enumerate}[nosep]
  \item build an MBRA runtime wrapper,
  \item prepare MBRA inputs from live EarthRover frames,
  \item run MBRA inference on recent context plus the subgoal image,
  \item convert MBRA outputs into EarthRover commands,
  \item clip and smooth commands before safety checks.
\end{enumerate}

\subsection{Exact role of MBRA}

\begin{mdframed}
\textbf{MBRA is the local controller.} It solves the short-horizon problem:
``from what I see now, how do I move toward this nearby goal image?''
\end{mdframed}

\subsection{Definition of done}

\begin{itemize}[nosep]
  \item robot can be driven toward a manually selected nearby indoor goal image,
  \item MBRA commands are correctly transformed into rover control commands,
  \item local goal-following is stable enough for integration with the planner.
\end{itemize}


\section{Stack 5: Safety}

\subsection{Purpose}

Prevent local control from producing unsafe motion.

\subsection{Inputs}

\begin{itemize}[nosep]
  \item current frame,
  \item MBRA proposed command,
  \item pedestrian detection flags from perception.
\end{itemize}

\subsection{Outputs}

\begin{itemize}[nosep]
  \item safe final linear command,
  \item safe final angular command,
  \item stop / slow / override reason.
\end{itemize}

\subsection{What we are going to do}

\begin{enumerate}[nosep]
  \item add a depth or clearance estimation step,
  \item define a forward safety check for commanded motion,
  \item stop or slow down when people are detected in the danger zone,
  \item provide a final safe command to the integration layer.
\end{enumerate}

\subsection{Preferred first safety rules}

\begin{itemize}[nosep]
  \item if a person is detected in the forward danger zone, slow or stop,
  \item if predicted clearance in commanded direction is too low, stop or redirect,
  \item if confidence is low or sensor state is stale, fail safe.
\end{itemize}

\subsection{Definition of done}

\begin{itemize}[nosep]
  \item unsafe MBRA commands are prevented,
  \item pedestrian-aware slowing or stopping works in corridor tests.
\end{itemize}


\section{Stack 6: Integration / Orchestration}

\subsection{Purpose}

Connect all stacks into one real system.

\subsection{Inputs}

\begin{itemize}[nosep]
  \item outputs from all other stacks.
\end{itemize}

\subsection{Outputs}

\begin{itemize}[nosep]
  \item one running indoor navigation process,
  \item logs,
  \item debug traces,
  \item replayable experiment data.
\end{itemize}

\subsection{What we are going to do}

\begin{enumerate}[nosep]
  \item define clear module interfaces,
  \item build one runtime loop that calls each stack in order,
  \item standardize config files and logging,
  \item provide debug views for localization, path, MBRA output, and safety.
\end{enumerate}

\subsection{Definition of done}

\begin{itemize}[nosep]
  \item all stacks can run together in one loop,
  \item end-to-end indoor tests are possible.
\end{itemize}


\section{Stack 7: Testing / Evaluation}

\subsection{Purpose}

Turn a collection of modules into a reliable system.

\subsection{What we are going to measure}

\begin{itemize}[nosep]
  \item localization accuracy against node labels,
  \item subgoal switching correctness,
  \item local goal-following success rate,
  \item safety interventions,
  \item checkpoint reach success rate,
  \item number and type of failures.
\end{itemize}

\subsection{What we are going to do}

\begin{enumerate}[nosep]
  \item define a few standard indoor routes,
  \item run repeated tests on those routes,
  \item log failures by stack,
  \item decide whether failures are from perception, planning, control, or safety.
\end{enumerate}


\section{Exact Handoffs Between Stacks}

\begin{longtable}{p{0.22\textwidth} p{0.30\textwidth} p{0.38\textwidth}}
\toprule
\textbf{From} & \textbf{To} & \textbf{Handoff} \\
\midrule
SDK & Perception & latest frame, timestamps, live stream access \\
Perception & Planning & current node estimate, candidate nodes, confidence \\
Map / Graph & Planning & graph structure, node images, target node ids \\
Planning & Control & current subgoal node id and subgoal image \\
Perception & Safety & pedestrian detections, checkpoint verification flags if relevant \\
Control & Safety & proposed rover command \\
Safety & Integration & final safe rover command \\
Integration & SDK & command send request \\
\bottomrule
\end{longtable}


\section{Immediate Engineering Order}

This is the order we should build in.

\begin{enumerate}
  \item \textbf{SDK stability} --- confirm live frame and control path are solid.
  \item \textbf{MBRA local prototype} --- prove the robot can move toward one manually selected nearby indoor goal image.
  \item \textbf{Graph build tooling} --- create one short indoor graph from collected data.
  \item \textbf{Visual localization} --- localize live frames into that graph.
  \item \textbf{Planner} --- choose subgoal nodes from current node to target node.
  \item \textbf{Safety layer} --- stop or slow on people and low clearance.
  \item \textbf{Full integration} --- run the whole loop.
  \item \textbf{Repeat testing} --- measure failure modes and refine.
\end{enumerate}


\section{What Each Team Member Should Produce}

\subsection*{SDK / Interface owner}
\begin{itemize}[nosep]
  \item stable camera and control API,
  \item frame logging utility,
  \item command logging utility.
\end{itemize}

\subsection*{Perception owner}
\begin{itemize}[nosep]
  \item node retrieval module,
  \item checkpoint verification module,
  \item pedestrian detector wrapper.
\end{itemize}

\subsection*{Map / Graph owner}
\begin{itemize}[nosep]
  \item graph schema,
  \item graph builder from recorded routes,
  \item graph serialization format.
\end{itemize}

\subsection*{Planning owner}
\begin{itemize}[nosep]
  \item graph search module,
  \item subgoal selection module,
  \item path progress tracker.
\end{itemize}

\subsection*{Control owner}
\begin{itemize}[nosep]
  \item MBRA wrapper,
  \item recent-frame buffer interface,
  \item MBRA-to-rover command conversion.
\end{itemize}

\subsection*{Safety owner}
\begin{itemize}[nosep]
  \item depth / clearance check,
  \item person-aware slow / stop logic,
  \item final command override policy.
\end{itemize}

\subsection*{Integration owner}
\begin{itemize}[nosep]
  \item one main runtime loop,
  \item stack wiring,
  \item unified configuration and logging.
\end{itemize}

\subsection*{Testing owner}
\begin{itemize}[nosep]
  \item test route definitions,
  \item evaluation metrics,
  \item failure reports and experiment tracking.
\end{itemize}


\section{What We Are Doing Exactly}

\begin{mdframed}
\textbf{Perception} tells us where we are and whether a checkpoint is reached.\\
\textbf{Map / Graph} stores the building as connected places.\\
\textbf{Planning} decides which place we should go to next.\\
\textbf{Control} uses MBRA to move toward that nearby place image.\\
\textbf{Safety} prevents bad local actions.\\
\textbf{Integration} makes everything run as one system.\\
\textbf{Testing} tells us whether any of this is actually reliable.
\end{mdframed}

\begin{mdframed}
\textbf{The single most important architectural rule:}
MBRA is our \textbf{local controller}, not our global planner.
\end{mdframed}

If all team members build their stack with that rule in mind, the project
stays coherent.

\end{document}
